\documentclass[11pt,a4paper]{article}
\usepackage[T1]{fontenc}
\usepackage[utf8]{inputenc}
\usepackage[spanish]{babel}
\usepackage{amsmath,amssymb}
\usepackage[a4paper,margin=2.3cm]{geometry}
\usepackage{graphicx}
\usepackage{tikz}
\usetikzlibrary{arrows.meta,babel}
\IfFileExists{cuna_60_visualizacion.pdf}{%
  \newcommand{\includevisualizacion}[1][]{\includegraphics[##1]{cuna_60_visualizacion.pdf}}%
}{%
  \usepackage[inkscapelatex=false]{svg}
  \newcommand{\includevisualizacion}[1][]{\includesvg[##1]{cuna_60_visualizacion}}%
}
\spanishdecimal{.}
\setlength{\parindent}{0pt}
\setlength{\parskip}{0.35em}
\setlength{\abovedisplayskip}{4pt}
\setlength{\belowdisplayskip}{4pt}
\setlength{\abovedisplayshortskip}{4pt}
\setlength{\belowdisplayshortskip}{4pt}

\newcommand{\encabezado}[3]{%
  {\Large\bfseries #1\par}
  \vspace{2pt}
  {\normalsize #2\hfill #3\par}
  \vspace{6pt}\hrule\vspace{8pt}}

\newcommand{\problema}[2]{%
  \vspace{0.55em}
  \noindent\textbf{#1.}\enspace{\small #2}\par
  \vspace{0.3em}}

\begin{document}
\encabezado{Potencial y densidad inducida en una cuña conductora de $60^\circ$}
{Electrodinámica I}{P.S., J.A., J.R.}

\problema{1}{Dos semiplanos conductores ideales conectados a tierra se intersectan sobre el
eje $z$ y delimitan la región $0<\phi<\alpha$, con $\alpha=\pi/3$. Una carga
puntual $q>0$ se encuentra en $(s,\phi,z)=(d,\pi/6,0)$, sobre la bisectriz.
Determine las cargas imagen, el potencial dentro de la cuña y la densidad
superficial inducida en ambas caras. Compruebe el comportamiento cerca de la
arista y lejos de la carga.}

\textbf{Solución:}

\textbf{1. Geometría y cargas imagen.}\par
Para $\alpha=\pi/N$ con $N=3$, las reflexiones sucesivas cierran un conjunto
de $2N=6$ cargas: una real y cinco imágenes. Es útil separarlas en dos
familias, para $k=0,1,2$:
\[
\theta_k^+=\phi_0+2k\alpha,\qquad q_k^+=+q;
\qquad
\theta_k^-=-\phi_0+2k\alpha,\qquad q_k^-=-q,
\quad \phi_0=\frac{\pi}{6}.
\]
\[
\begin{array}{c|ccc}
 & k=0 & k=1 & k=2\\ \hline
\theta_k^+ & 30^\circ\;(\text{real}) & 150^\circ & 270^\circ\\
\theta_k^- & 330^\circ & 90^\circ & 210^\circ
\end{array}
\]
Todas se sitúan en $s=d$, $z=0$. Solo la carga en $30^\circ$ pertenece a la
región física.

\begin{figure}[htbp]
\centering
\begin{tikzpicture}[
  scale=1.05,>=Stealth,
  line cap=round,line join=round,
  every node/.style={font=\small}
]
  \def\dist{2.02}
  % Solo el sector sombreado es la región física.
  \fill[black!8] (0,0)--(0:2.67) arc (0:60:2.67)--cycle;
  \draw[->,thin] (-2.75,0)--(3.03,0) node[right] {$x$};
  \draw[->,thin] (0,-2.75)--(0,2.98) node[above] {$y$};
  \draw[densely dashed,black!55] (0,0) circle (\dist);
  \draw[very thick] (0,0)--(0:2.82)
       node[pos=0.96,below] {$\phi=0$};
  \draw[very thick] (0,0)--(60:2.82)
       node[above right] {$\phi=\alpha$};
  \draw[->] (0.62,0) arc (0:60:0.62);
  \node at (28:0.87) {$60^\circ$};
  \draw[densely dotted] (0,0)--(30:\dist)
       node[pos=0.80,below] {$d$};

  % Las flechas apuntan desde cada conductor hacia el vacío.
  \draw[->,semithick] (2.25,0)--++(90:0.39)
       node[above] {$\hat{\mathbf n}_0$};
  \draw[->,semithick] (60:2.25)--++(-30:0.39)
       node[below right] {$\hat{\mathbf n}_{\alpha}$};

  % Punto lleno: carga real. Puntos vacíos: cargas imagen.
  \foreach \a in {150,270}{
    \draw[thick,fill=white] (\a:\dist) circle (2.6pt);
    \node[fill=white,inner sep=1pt] at (\a:2.43) {$+q$};
  }
  \foreach \a in {90,210,330}{
    \draw[thick,fill=white] (\a:\dist) circle (2.6pt);
    \node[fill=white,inner sep=1pt] at (\a:2.43) {$-q$};
  }
  \fill (30:\dist) circle (2.8pt);
  \node[anchor=west,fill=white,inner sep=1pt] at (2.45,1.05)
       {$+q$ real};
  \fill (0,0) circle (1.2pt) node[below left] {$O$};
\end{tikzpicture}
\caption{Corte perpendicular a la arista $z$. El sector sombreado es la
región física $0<\phi<\alpha$; el punto negro es la carga real y los
puntos blancos son imágenes. La circunferencia discontinua solo indica
$s=d$. Las flechas señalan las normales de los conductores hacia la cuña.}
\label{fig:imagenes-cuna}
\end{figure}

\textbf{2. Potencial y condiciones de borde.}\par
La distancia entre $(s,\phi,z)$ y una carga en $(d,\theta,0)$ es
\[
R(s,\phi,z;\theta)=
\sqrt{s^2+d^2-2sd\cos(\phi-\theta)+z^2}.
\]
La solución dentro de $0<\phi<\pi/3$ es
\begin{equation}
\boxed{\displaystyle
V(s,\phi,z)=\frac{q}{4\pi\varepsilon_0}
\sum_{k=0}^{2}
\left[\frac{1}{R(s,\phi,z;\theta_k^+)}
     -\frac{1}{R(s,\phi,z;\theta_k^-)}\right].}
\label{eq:potential}
\end{equation}
En $\phi=0$ cada carga en $\theta$ se empareja con otra de signo opuesto en
$-\theta$, a igual distancia del punto de observación. La reflexión respecto
de $\phi=\alpha$ da el mismo resultado para la otra pared. Por lo tanto,
$V(s,0,z)=V(s,\alpha,z)=0$. La solución cumple la ecuación de Poisson en la
carga real, Laplace en el resto de la cuña y decae en el infinito.

\textbf{3. Densidad inducida: normales y signo.}\par
La derivada correcta de cada término es
\[
\frac{\partial}{\partial\phi}\frac{1}{R(s,\phi,z;\theta)}
=-\frac{sd\sin(\phi-\theta)}{R(s,\phi,z;\theta)^3},
\qquad
E_\phi=-\frac{1}{s}\frac{\partial V}{\partial\phi}.
\]
La normal que sale del conductor hacia la cuña es $+\hat{\boldsymbol\phi}$
en $\phi=0$ y $-\hat{\boldsymbol\phi}$ en $\phi=\alpha$. Así,
\[
\sigma_0(s,z)=\varepsilon_0 E_\phi(s,0,z),\qquad
\sigma_\alpha(s,z)=-\varepsilon_0 E_\phi(s,\alpha,z).
\]
Definamos
\[
C=s^2+d^2+z^2,\qquad
A=C-\sqrt{3}\,sd,\qquad B=C+\sqrt{3}\,sd.
\]
Al sustituir las seis cargas y agrupar pares reflejados se obtiene
\begin{equation}
\boxed{\displaystyle
\sigma_0(s,z)=\sigma_\alpha(s,z)=
-\frac{qd}{4\pi}
\left(A^{-3/2}+B^{-3/2}-2C^{-3/2}\right).}
\label{eq:sigma}
\end{equation}
La igualdad de las dos caras también sigue de la simetría respecto de la
bisectriz. Como $u^{-3/2}$ es estrictamente convexa para $u>0$ y
$C=(A+B)/2$, la expresión entre paréntesis es positiva para $s>0$:
ambas paredes adquieren carga inducida negativa si $q>0$.

\textbf{4. Chequeos físicos.}\par
En la arista, con $z$ fijo y $s\to0^+$, la expansión de
\eqref{eq:sigma} da
\[
\sigma(s,z)=
-\frac{45qd^3}{16\pi(d^2+z^2)^{7/2}}\,s^2+O(s^4).
\]
La densidad \emph{tiende a cero}, no diverge, en esta cuña de $60^\circ$.
En cambio, para $s\to\infty$ a $z$ fijo,
$\sigma(s,z)\sim-45qd^3/(16\pi s^5)$.
Las imágenes cancelan los momentos lejanos de orden bajo, de modo que el
flujo eléctrico a través de una superficie lejana se anula. Por la ley de
Gauss, la carga total inducida entre ambas caras es $-q$; la simetría
asigna $-q/2$ a cada cara.

\clearpage
\section*{Visualización numérica}
El script Python usa $d=q=4\pi\varepsilon_0=1$ para evaluar
\eqref{eq:potential} y \eqref{eq:sigma}. Las cargas fuera de la cuña en el
panel geométrico son imágenes matemáticas. En el mapa del potencial se
excluye un disco pequeño alrededor de la carga real, donde el modelo puntual
es singular. Los dos mapas y los perfiles se calculan a partir de las mismas
fórmulas analíticas.
\begin{figure}[htbp]
\centering
\includevisualizacion[width=0.98\textwidth]
\caption{(a) Posición y signo de las seis cargas. (b) Potencial normalizado
dentro de la cuña, nulo en ambas paredes. (c) Densidad inducida sobre
$z=0$ en función de la distancia a la arista; las dos caras coinciden.
(d) Densidad inducida sobre una cara en función de $s$ y $z$.}
\label{fig:numerica}
\end{figure}
\end{document}
